% Generated draft: requires mathematical and source review.
\documentclass[11pt]{article}
\usepackage[margin=25mm]{geometry}
\usepackage{fontspec}
\setmainfont{Latin Modern Roman}
\usepackage{amsmath,amssymb,mathtools,unicode-math}
\setmathfont{Latin Modern Math}
\usepackage{enumitem,needspace,xurl,hyperref}
\setlength{\parindent}{0pt}
\setlength{\parskip}{5pt}
\setlength{\emergencystretch}{3em}
\newcommand{\N}{\mathbb{N}}
\newcommand{\Z}{\mathbb{Z}}
\newcommand{\Q}{\mathbb{Q}}
\newcommand{\R}{\mathbb{R}}
\newcommand{\C}{\mathbb{C}}
\newcommand{\F}{\mathbb{F}}
\newcommand{\A}{\mathbb{A}}
\newcommand{\PP}{\mathbb{P}}
\newcommand{\E}{\mathbb{E}}
\newcommand{\eps}{\varepsilon}
\newcommand{\dd}{\,\mathrm{d}}
\begin{document}
\section*{UnsolvedMath ID: 30000433}
\textbf{Source problem code:} catalogue entry
\subsection{Definitions and mathematical statement}
\textbf{Two triangulation questions.} Does every three-manifold admit a triangulation in which every edge has degree \(5\) or \(6\)? Can such a triangulation moreover be chosen so that no triangle contains more than one degree-\(6\) edge? The second is a stronger optional assertion, not a replacement for the first.\par For a simplicial triangulation \(K\) with \(|K|\cong M\), define \(\deg_K(e)=\#\{\sigma\in K:\dim\sigma=3,\ e\subset\sigma\}\). Thus degree counts incident tetrahedra, not the number of endpoints, neighboring vertices, or a metric dihedral angle. The basic requirement is\[\deg_K(e)\in\{5,6\}\quad\text{for every edge }e.\] The stronger condition is\[\#\{e\subset\tau:\dim e=1,\ \deg_K(e)=6\}\le1\quad\text{for every triangle }\tau\in K.\] These are combinatorial conditions; no regular-tetrahedron metric or geometric embedding is imposed.\par The source says every three-manifold but does not fix closed versus boundary cases, compact versus noncompact cases, orientability, or genuine simplicial complexes versus tetrahedral face-pairing conventions. For compact \(M\) one convention uses finite triangulations; noncompact \(M\) requires a local-finiteness convention. In a generalized face-pairing triangulation, incidence multiplicities must be counted explicitly, whereas the formula above uses a genuine simplicial complex. The intended scope must be checked rather than silently restricted to a favored subclass.\par The supplied literature triage checked 2026-08-21 labels both universal assertions open. It reports degree-\(4,5,6\) triangulations in general settings and many examples satisfying the stronger condition, called TCP examples in that assessment. Exact hypotheses of the general construction and the TCP terminology are not reproduced. The reported \(\{4,5,6\}\) result does not by itself remove degree \(4\), nor do examples establish the all-manifolds claim. This dated assessment requires independent verification.
\subsection{Short English statement}
Ask first for a triangulation with five or six tetrahedra at every edge, and then for the stronger restriction of at most one degree-six edge in each triangle. The manifold and triangulation conventions are unspecified; no geometric regularity of tetrahedra is assumed.
\subsection{Sources}
\begin{itemize}\item Local UnsolvedMath registry, permanent ID 30000433, code catalogue entry; \emph{Five-or-Six Edge-Degree Triangulations of Three-Manifolds}, Oberwolfach Reports Open Problems.\item \emph{Discrete Differential Geometry (2006)}, Oberwolfach Reports, DOI 10.4171/owr/2006/12, pp. 653--728. \url{https://doi.org/10.4171/owr/2006/12}. Proposed year 2006; no proposer named.\item Supplied literature label \emph{original OWR problem}: \url{https://doi.org/10.4171/owr/2006/12}.\item Supplied literature label \emph{Brady--McCammond--Meier edge degrees}: \url{https://math.ou.edu/~nbrady/papers/edgedegrees.pdf}.\item Literature triage checked 2026-08-21, status \texttt{open}; record updated 2026-09-16. Statement status \texttt{exact}, with source checking reported in the verification field, not independently checked here.\end{itemize}
\end{document}
